% Source-derived chapter 02: A uniformity estimate for odd degree monic polynomials
% Original source: squarefree-values-polynomial-discriminants
\chapter{A uniformity estimate for odd degree monic polynomials}
\label{squarefree-values-polynomial-discriminants-chapter-02}
\source{squarefree-values-polynomial-discriminants}

\begin{remark}[Source section]
\label{squarefree-values-polynomial-discriminants-chapter-02-source}
\source{squarefree-values-polynomial-discriminants}
\source{squarefree-values-polynomial-discriminants}
This chapter reproduces the corresponding section of the retrieved
LaTeX source, including its definitions, statements, examples, and
proofs. It has not yet been translated into Lean.
\notready
\end{remark}

% The source section title was: A uniformity estimate for odd degree monic polynomials
\label{sec:monicodd}
\source{squarefree-values-polynomial-discriminants}

In this section, we prove the estimate of
Theorem~\ref{thm:mainestimate}(b) when $n=2g+1$ is odd, for any $g\geq
1$.

\subsection{Invariant theory for the fundamental representation:
  $\SO_n$ on the space $W$ of symmetric $n\times n$ matrices}

Let $A_0$ denote the $n\times n$ symmetric matrix with $1$'s on the
anti-diagonal and $0$'s elsewhere. The group $G=\SO(A_0)$ acts on $W$
via the action
\begin{equation*}
\gamma\cdot B=\gamma B\gamma^t.
\end{equation*}

We recall some of the arithmetic invariant theory for the
representation $W$ of $n\times n$ symmetric matrices of the split
orthogonal group $G$; see \cite{BG2} for more details.  The ring of
polynomial invariants for the action of $G(\C)$ on $W(\C)$ is freely
generated by the coefficients of the {\it invariant polynomial
  $f_B(x)$ of $B$}, defined by $$f_B(x):=(-1)^{g}\det(A_0x-B)$$ (see
\cite[\S4]{AITBG}). We define the {\it discriminant} $\Delta$ on $W$
by $\Delta(B)=\Delta(f_B)$, and the $G(\R)$-invariant {\it height} of
elements in $W(\R)$ by $H(B)=H(f_B).$


Let $k$ be any field of characteristic not $2$.
%Then elements $B\in W(k)$
%with $\Delta(B)\neq0$ are stable in the sense of geometric invariant
%theory.
For a monic polynomial $f(x)\in k[x]$ of degree~$n$ such that
$\Delta(f)\neq0$, let $C_f$ denote the smooth hyperelliptic curve
$y^2=f(x)$ of genus $g$ and let $J_f$ denote the Jacobian of
$C_f$. Then $C_f$ has a rational Weierstrass point at infinity.
%The set of $G(k)$-orbits in $W(k)$ having invariant polynomial $f(x)$
%is a subset of $H^1(k,J_f)$.
The stabilizer of an element $B\in W(k)$ with invariant polynomial
$f(x)$ is naturally isomorphic to $J_f[2](k)$
by~\cite[Proposition~5.1]{BG2}, and hence has cardinality at most
$\#J_f[2](\bar k)=2^{2g}$, where $\bar k$ denotes a separable closure
of $k$.

We say that an element (or the $G(k)$-orbit of an element) $B\in W(k)$ with $\Delta(B)\neq 0$
is {\it $k$-distinguished} if there exists a $g$-dimensional subspace
defined over $k$ that is isotropic with respect to both $A_0$ and~$B$. If $B$ is $k$-distinguished, then the set of these $g$-dimensional
subspaces over $k$ is in bijection with $J_f[2](k)$
by~\cite[Proposition~4.1]{BG2}, and so it too has cardinality at most
$2^{2g}$.

In fact, it is known (see \cite[Proposition~5.1]{BG2}) that the
elements of $J_f[2](k)$ are in natural bijection with the even-degree
factors of $f$ defined over $k$.  (Note that the number of even-degree
factors of $f$ over $\bar k$ is indeed $2^{2g}$.) In particular, if
$f$ is irreducible over $k$, then the group $J_f[2](k)$ is trivial.

Now let $W_0$ be the subspace of $W$ consisting of matrices whose top left
$g\times g$ block is zero. Then elements $B$ in $W_0(k)$ with nonzero
discriminant are all evidently $k$-distinguished since the $g$-dimensional subspace
$Y_g$ spanned by the first $g$ basis vectors is isotropic with respect
to both $A_0$ and $B$.
Let $G_0$ denote the subgroup of $G$
consisting of elements $\gamma$ such that $\gamma^t$ preserves
$Y_g$. Then $G_0$ acts on $W_0$.

An element $\gamma\in G_0$ has the block matrix form
\begin{equation}\label{eq:G_0}
\source{squarefree-values-polynomial-discriminants}
\gamma=\Bigl(\begin{array}{cc}\gamma_1 & 0\\ \delta & \gamma_2
\end{array}\Bigr)\in\Bigl(\begin{array}{cc}M_{g\times g} & 0\\ M_{(g+1)\times g} & M_{(g+1)\times (g+1)}
\end{array}\Bigr),
\end{equation}
and so $\gamma\in G_0$ transforms the top right $g\times (g+1)$ block
of an element $B\in W_0$ as follows:
$$(\gamma\cdot B)^\sub = \gamma_1B^\sub\gamma_2^t,$$ where we use the
superscript ``top'' to denote the top right $g\times (g+1)$ block of
any given element in $W_0$.  It will be convenient for us to view
$(A_0^\sub,B^\sub)$ as an element of the representation $V_g=2\otimes
g\otimes (g+1)$ of the group $H_g:=\SL_2\times\GL_g\times\GL_{g+1}$.
We have a map $\theta:G_0\to H_g$ sending $\gamma$ expressed in \eqref{eq:G_0} to
$(1,\gamma_1,\gamma_2)$. Then we have $$(A_0^\sub,(\gamma\cdot
B)^\sub)=\theta(\gamma)\cdot(A_0^\sub,B^\sub)$$ for $\gamma\in G_0$ and
$B\in W_0$.

%Then the action of $\gamma\in G_0$ on $B\in
%W_0$ agrees with the action of $\theta(\gamma)$ on the pair
%$(A_0^\sub,B^\sub)$.

Next, we construct a relative polynomial invariant for the action of $H_g$ on
$V_g$ as follows.  We write any $2\times g\times (g+1)$ matrix $v$ in
$V_g$ as a pair $(A,B)$ of $g\times(g+1)$ matrices.  Let $M_v(x,y)$
denote the vector of $g\times g$ minors of $Ax-By$, where $x$ and $y$
are indeterminates; in other words, the $i$-th coordinate of the
vector $M_v(x,y)$ is given by $(-1)^{i-1}$ times the determinant of
the matrix obtained by removing the $i$-th column of $Ax-By$.  Then
$M_v(x,y)$ is a vector of length $g+1$ consisting of binary forms of
degree $g$ in $x$ and~$y$, each of which has $g+1$ coefficients.
Taking the determinant of the resulting $(g+1)\times(g+1)$ matrix of
coefficients of these $g+1$ binary forms in $M_v(x,y)$ then yields a
polynomial $Q=Q(v)$ in the coordinates of $V_g$, which is a relative
invariant for the action of $H_g$. Explicitly, we have
\begin{equation}\label{eq:relQ}
\source{squarefree-values-polynomial-discriminants}
Q((\gamma_0,\gamma_1,\gamma_2)\cdot v)=\det(\gamma_1)^{g+1}\det(\gamma_2)^g Q(v).
\end{equation}
Indeed, it follows from the definition that the $Q$-invariant is
invariant under the action of the subgroup
$\SL_2\times\SL_g\times\SL_{g+1}$. (Alternatively, the group
$\SL_2\times\SL_g\times\SL_{g+1}$ has no nontrivial characters.)
Moreover, we may work over the algebraic closure and write $\gamma_1$
and $\gamma_2$ as products of scalar matrices and matrices with
determinant 1.  Finally one can easily check \eqref{eq:relQ} when
$\gamma_1$ and $\gamma_2$ are scalar matrices.
%We call this polynomial the $Q$-{\it invariant}.

We then define the $Q$-{\it invariant} of $B\in W_0$ to be the
$Q$-invariant of $(A_0^\sub,B^\sub)$:
\begin{equation}\label{eqQB}
\source{squarefree-values-polynomial-discriminants}
Q(B):=Q(A_0^\sub,B^\sub).
\end{equation}
Then the $Q$-invariant is also a relative invariant for the action of $G_0$ on $W_0$, since for any
$\gamma\in G_0$ expressed in the form \eqref{eq:G_0}, we have
\begin{equation}\label{eq:weightG_0}
\source{squarefree-values-polynomial-discriminants}
Q(\gamma\cdot B) = \det(\gamma_1)Q(B).
\end{equation}
In fact, we may extend the definition of the $Q$-invariant to an even
larger subset of $W(\Q)$ than $W_0(\Q)$. We have the following
proposition.

\begin{proposition}
\source{squarefree-values-polynomial-discriminants}
\notready\label{prop:extendQ}
\source{squarefree-values-polynomial-discriminants}
Let $B\in W_0(\Q)$ be an element whose invariant polynomial $f(x)$ is
irreducible over~$\Q$. Then for every $B'\in W_0(\Q)$ such that $B'$ is $G(\Z)$-equivalent to $B$, we have
$Q(B')=\pm Q(B)$.
\end{proposition}
\begin{proof}
Suppose $B'=\gamma\cdot B$ with $\gamma\in G(\Z)$ and $B,B'\in
W_0(\Q)$.  Then $Y_g$ and $\gamma^t Y_g$ are both $g$-dimensional
subspaces over $\Q$ isotropic with respect to both $A_0$ and
$B$. Since $f$ is irreducible over $\Q$, we have that $J_f[2](\Q)$ is
trivial, and so these two subspaces must be the same.  We conclude
that $\gamma\in G_0(\Z)$, and thus $Q(\gamma\cdot B)=\pm Q(B)$
by~\eqref{eq:weightG_0}.
\end{proof}

We may thus define the $|Q|$-{\it invariant} for any element $B\in
W(\Q)$ that is $G(\Z)$-equivalent to some element $B'\in W_0(\Q)$ and
whose invariant polynomial is irreducible over $\Q$; indeed, we set
$|Q|(B):=|Q(B')|$. By Proposition~\ref{prop:extendQ}, this definition
of $|Q|(B)$ is independent of the choice of $B'$. Note that all such
elements $B\in W(\Q)$ are $\Q$-distinguished.

\subsection{Embedding $\U_m^\w$ into $\frac12W(\Z)$}\label{sembedodd}
\source{squarefree-values-polynomial-discriminants}

We begin by describing those monic integer polynomials in $\Vmn(\Z)$
that lie in $\U_m^\w$, i.e., the monic integer polynomials that have
discriminant weakly divisible by $p^2$ for all $p\mid m$.

\begin{proposition}
\source{squarefree-values-polynomial-discriminants}
\notready\label{prop:hiddenp2}
\source{squarefree-values-polynomial-discriminants}
Let $m$ be a positive squarefree integer, and let $f$ be a monic
integer polynomial whose discriminant is weakly divisible by $p^2$ for all $p\mid m$.
Then there exists an integer $\ell$ such that $f(x+\ell )$ has
the form
\begin{equation}
f(x+\ell ) = x^n + b_1x^{n-1} + \cdots + b_{n-2}x^2 + b_{n-1}x +
b_n
\end{equation}
for some integers $b_1,\ldots,b_n$ where $m$ divides $b_{n-1}$ and $m^2$ divides $b_n$.
\end{proposition}
\begin{proof}
Since $m$ is squarefree, by the Chinese Remainder Theorem it suffices
to prove the assertion in the case that $m=p$ is prime. Since $p$
divides the discriminant of~$f$, the reduction of $f$ modulo~$p$ must
have a repeated factor $h(x)^e$ for some polynomial $h\in \F_p[x]$ and
some integer $e\geq2$. As the discriminant of $f$ is not strongly
divisible by $p^2$, we see that $h$ is linear and $e=2$. By replacing
$f(x)$ by $f(x+\ell )$ for some integer $\ell$, if necessary, we may
assume that the repeated factor is $x^2$, i.e., we may assume that
$f(x)$ has the form $$f(x) = x^n + b_1x^{n-1} + \cdots + b_{n-1}x +
b_n$$ for some integers $b_1,\ldots,b_n$ such that $p$ divides
$b_{n-1}$ and $b_n$. It remains now to show that $p^2$ divides~$b_n$.

Viewing the discriminant $\Delta(f)$ has a polynomial in $b_n$,
we write
$$\Delta(f) = b_n\Delta_1 + \Delta_2,$$ for polynomials $\Delta_1\in
\Z[b_1,\ldots,b_n]$ and $\Delta_2\in \Z[b_1,\ldots,b_{n-1}]$. Next we
set $b_n$ to $0$ and observe that $b_{n-1}^2$ divides
$\Delta_2$. Indeed, the discriminant of $f$ is equal to the resultant
$\Res(f(x),f'(x))$ of $f(x)$ and $f'(x)$.  When $b_n=0$, the only
nonzero entry in the last column of the matrix whose determinant
computes $\Res(f(x),f'(x))$ is $b_{n-1}$ appearing in the last
row. The corresponding minor after removing the last row and column
has two nonzero entries in its last column, both of which equal to
$b_{n-1}$. This implies that we can pull out another factor of $b_{n-1}$ from the determinant of this minor.
Hence there is a polynomial $\Delta_3\in
\Z[b_1,\ldots,b_{n-1}]$ such that
\begin{equation}\Delta(f) = b_n\Delta_1 + b_{n-1}^2\Delta_3.
\end{equation}
Since $p$ divides $b_{n-1}$ and $b_n$, we see that $p^2\mid \Delta(f)$
if and only if $p^2\mid b_n\Delta_1$. If $p^2$ does not divide $b_n$,
then $p$ divides $\Delta_1$, which implies that $p^2$ divides $\Delta(f)$
strongly, a contradiction. Therefore, we have that $p^2$ divides $b_n$.
\end{proof}

Proposition \ref{prop:hiddenp2} identifies the $p^2$ hidden in the
coefficients of a monic polynomial $f$ when $p^2$ weakly divides
$\Delta(f)$. We next construct a matrix in $\frac12 W_0(\Z)$ with
invariant polynomial $f$ and where the two $p$'s in $p^2$ are split
apart.  For any integers $m,c_1,\ldots,c_n$, consider the matrix
\begin{equation}\label{mat1}
\source{squarefree-values-polynomial-discriminants}
 B_m(c_1,\ldots,c_n) =
 \left(\!\begin{array}{ccccccccc}&&&&&&&m&\!0\\[.1in]&&&&&&\iddots&\iddots& \\[.15in]&&&&&\,1\,&\,0\,&&\\[.125in] &&&&\,1\,&0&&&\\[.125in] &&&1&-c_1&\!\!-c_2/2\!\!&&& \\[.15in] &&1&\;\,0\;\,&\!\!-c_2/2\!\!&-c_3&\!\!-c_4/2\!\!&&\\[.045in]&\iddots&\;\,\,0\;\,\,&&&\!\!-c_4/2\!\!&-c_5&\ddots&\\[.025in] \;m\;&\,\,\,\iddots\,\,\,&&&&&\ddots&\ddots&\!\!\!-c_{n-1}/2\!\!\!
   \\[.105in] \,0\,&&&&&&&\!\!\!-c_{n-1}/2\!\!\!&-c_n \end{array}\right)
\end{equation}
in $\frac12 W_0(\Z)$. It follows from a direct computation that
$$f_{B_m(c_1,\ldots,c_n)}(x) = x^n + c_1x^{n-1} + \cdots + c_{n-2}x^2
+ mc_{n-1}x + m^2c_n.$$ We compute the $Q$-invariant of $B_m(c_1,\ldots,c_n)$. Let $v = (A_0^\sub, B_m(c_1,\ldots,c_n)^\sub)$ be the corresponding element in $V_g$. The vector $M_v(x,y)$ of length $g+1$ of minors of $A_0^\sub x-B_m(c_1,\ldots,c_n)^\sub y$ is
$$M_v(x,y) = \begin{pmatrix} x^g\\ x^{g-1}y\\ \vdots \\ xy^{g-1}\\ my^g\end{pmatrix}.$$
Computing the determinant of the $(g+1)\times(g+1)$ matrix of coefficients of $M_v(x,y)$ then gives
$$Q(B_m(c_1,\ldots,c_n)) = Q(v) = m.$$
For any integer $\ell$, we have
\begin{equation}\label{eq:ell}
\source{squarefree-values-polynomial-discriminants}
f_{B_m(c_1,\ldots,c_n) + \ell A_0}(x) = (-1)^g\det(A_0(x-\ell)-B_m(c_1,\ldots,c_n)) = f_{B_m(c_1,\ldots,c_n)}(x - \ell),
\end{equation}
and by the $\SL_2$-invariance of the $Q$-invariant,
\begin{equation}\label{eq:ellQ}
\source{squarefree-values-polynomial-discriminants}
Q(B_m(c_1,\ldots,c_n) + \ell A_0) = Q(B_m(c_1,\ldots,c_n)) = m.
\end{equation}

\begin{theorem}
\source{squarefree-values-polynomial-discriminants}
\notready\label{keymap}
\source{squarefree-values-polynomial-discriminants}
Let $m$ be a positive squarefree integer. There exists a map
$\sigma_m:\U_m^\w\to \frac12W_0(\Z)$ such that for every $f\in \U_m^\w$,
\begin{equation}\label{eq:liftQ}
\source{squarefree-values-polynomial-discriminants}
f_{\sigma_m(f)}=f,\qquad Q(\sigma_m(f)) = m.
\end{equation}
\end{theorem}

\begin{proof}
Let $f$ be any element of $\U_m^\w$. By Proposition \ref{prop:hiddenp2}, there exists an integer $\ell$ and integers $c_1,\ldots,c_n$ such that
$$f(x+\ell) = x^n + c_1x^{n-1} + \cdots + c_{n-2}x^2 + mc_{n-1}x + m^2c_n.$$
We set $\sigma_m(f) = B_m(c_1,\ldots,c_n) + \ell A_0$. Then \eqref{eq:liftQ} follows from  \eqref{eq:ell} and \eqref{eq:ellQ}.
\end{proof}

We remark that the integer $\ell$ in the proof of Theorem \ref{keymap}
is unique modulo $m$ since it is the unique double root of $f(x)$
modulo every prime factor $p$ of $m$. Since the invariant polynomial
of $\sigma_m(f)$ recovers $f$, we see that no two elements in the
image of $\sigma_m$ are in the same $G(\Z)$-orbit (or even
$G(\C)$-orbit). Furthermore, by the definition of discriminant on $W$,
we have that $\Delta(\sigma_m(f))=\Delta(f)$ is divisible by $p^2$ for
every prime factor $p$ of $m$. If one varies the coefficients of
$\sigma_m(f)$ by multiples of $p$ while keeping the top left $g\times
g$ block $0$, one can show that $p^2$ still divides the discriminant
of the resulting matrix. In other words, the matrix $\sigma_m(f)$, as
an element of $\frac12 W_0(\Z)$, has discriminant \emph{strongly
  divisible} by $p^2$. It then seems natural to use the Ekedahl sieve
to handle this strongly divisible case. However, the Ekedahl sieve
does not apply when the polynomial is not squarefree as a polynomial
and as we will show in the sequel, as polynomials in the coordinates
of $W_0$, the discriminant $\Delta$ is divisible by $Q^2$. Instead, we
will count elements in $\frac12W_0(\Z)$ having large $Q$-invariant
using geometry-of-numbers techniques.

More precisely, let $\LL$ be the set of elements $v\in \frac12W(\Z)$
satisfying the following conditions: $v$ is $G(\Z)$-equivalent to some
element in $\frac12W_0(\Z)$ and the invariant polynomial of $v$ is
irreducible over $\Q$.  Then by the remark following
Proposition~\ref{prop:extendQ}, we may view $|Q|$ as a function also
on $\LL$.  Using $\U_m^{\w,\irr}$ to denote the set of irreducible
polynomials in $\U_m^\w$, we then have the following immediate
consequence of Theorem \ref{keymap}:

\begin{theorem}
\source{squarefree-values-polynomial-discriminants}
\notready\label{keymaporbit}
\source{squarefree-values-polynomial-discriminants}
Let $m$ be a positive squarefree integer. There exists an injective map
$$\bar{\sigma}_m:\U_m^{\w,\irr}\to G(\Z)\backslash\LL$$ such that
$f_{\bar{\sigma}_m(f)}=f$ for every $f\in \U_m^{\w,\irr}$. Moreover,
for every element $B$ in the $G(\Z)$-orbit of an element in the image
of $\bar{\sigma}_m$, we have $|Q|(B)=m$.
\end{theorem}

The number of reducible monic integer
polynomials having height less than $X$ is of a strictly smaller order
of magnitude than the total number of such polynomials (see, e.g.,
Proposition~\ref{propredboundall}).  Thus, for our purposes of proving
Theorem~\ref{thm:mainestimate}(b), it will suffice to count elements
in $\U_m^{\w,\irr}$ of height less than $X$ over all $m>M$, which by
Theorem~\ref{keymaporbit} we may do by counting these special
$G(\Z)$-orbits on $\LL\subset \frac12W(\Z)$ having height less than
$X$ and $|Q|$-invariant greater than $M$.  More precisely, let
$N(\LL;M;X)$ denote the number of $G(\Z)$-equivalence classes of
elements in $\LL$ whose $|Q|$-invariant is greater than $M$ and whose
height is less than $X$.  Then, by Theorem~\ref{keymaporbit}, to
obtain an upper bound for the left hand side in
Theorem~\ref{thm:mainestimate}(b), it suffices to obtain the same
upper bound for $N(\LL;M;X)$.

\subsection{Counting $G(\Z)$-orbits in $\frac12W(\Z)$}

The counting problem for the representation $W$ of $G$ is studied in
\cite{BG2}. In this section, we recall some of the set up and results
of \cite{BG2}.

We remark first that the representation studied in \cite{BG2} is the
subspace of $W$ consisting of symmetric matrices with anti-trace
$0$. Since this is a linear condition on the anti-diagonal entries,
each of which has weight $1$ when restricted to the action of the
maximal torus (see also \eqref{wbij}), the only difference in not
imposing this condition is an extra factor of $X$ in the count. We
will in fact make use of the anti-trace-$0$ version in Section~\ref{latticearg} in our
application to counting fields.

To count $G(\Z)$-orbits in a lattice $\frac12W(\Z)$ in $W(\R)$, one
begins by constructing fundamental domains for the action of $G(\Z)$
on the set of elements in $W(\R)$ with nonzero discriminant. A
fundamental domain $R$ for the action of $G(\R)$ on the set of
elements in $W(\R)$ with nonzero discriminant and height less than $1$
is obtained in \cite[\S9.1]{BG2}. The exact shape of $R$ is not
important. What is important is that $R$ is (absolutely) bounded. Next
a fundamental domain $\FF$ for the left multiplication action of
$G(\Z)$ on $G(\R)$ is written down in \cite[\S9.2]{BG2}. This is done
using the Iwasawa decomposition of $G(\R)$ as
$$
G(\R)=N(\R)TK,
$$ where $N$ is a unipotent group consisting of lower triangular matrices, $K$ is compact, and $T$ is the split torus of $G$ given by
\begin{equation*}
T=
\left\{\left(\begin{array}{ccccccc}
 t_1^{-1}&&&&&&\\
&\ddots &&&&&  \\
 && t_{g}^{-1} &&&&\\
&&& 1 &&&\\
 &&&& t_g &&\\
&&&&&\ddots &  \\
& &&&&& t_{1}
\end{array}\right):t_1,\ldots,t_g\in\R\right\}.
\end{equation*}
We may also make the following change of variables. For $1\leq
i\leq g-1$, set $s_i$ to be
$$
s_i=t_i/t_{i+1},
$$ and set $s_g=t_g$.  It follows that for $1\leq i\leq g$, we have
$t_i = s_is_{i+1}\cdots s_g.$ We denote an element of $T$ with
coordinates $t_i$ (resp.\ $s_i$) by $(t)$ (resp.\ $(s)$). A
fundamental set for the action of $G(\Z)$ on $G(\R)$ can then be taken
to be contained in a {\it Siegel set}, i.e., contained in $N'T'K$,
where $N'$ consists of elements in $N(\R)$ whose coefficients are
absolutely bounded and $T'\subset T$ consists of elements in $(s)\in
T$ with $s_i\geq c$ for some positive constant $c$.

For any $h\in G(\R)$, since $\FF h$ remains a fundamental domain for
the action of $G(\Z)$ on $G(\R)$, the set $(\FF h)\cdot (XR)$ (when
viewed as a multiset) is a finite cover of a fundamental domain for
the action of $G(\Z)$ on the elements in $W(\R)$ with nonzero
discriminant and height bounded by $X$. The degree of the cover
depends only on the size of stabilizer in $G(\R)$ and is thus
absolutely bounded by $2^{2g}$. The presence of these stabilizers is
in fact the reason we consider $\FF h\cdot XR$ as a multiset. Hence,
we have
\begin{equation}\label{eqoddfundv}
\source{squarefree-values-polynomial-discriminants}
  N(\LL;M;X)\ll \#\{B\in ((\FF h)\cdot (XR))\cap\LL:|Q|(B)>M\}.
\end{equation}
Let $\GG_1$ be a compact left $K$-invariant set in $G(\R)$ which is the
closure of a nonempty open set.
Averaging \eqref{eqoddfundv} over $h\in \GG_1$ and exchanging the
order of integration as in \cite[\S10.1]{BG2}, we obtain
\begin{equation}
  N(\LL;M;X)\ll \int_{\gamma\in\FF} \#\{B\in((\gamma \GG_1)\cdot (XR))\cap\LL:|Q|(B)>M\}
  d\gamma,
\end{equation}
where the implied constant depends only on $\GG_1$ and $R$, and where $d\gamma$ is a Haar measure on $G(\R)$ given by
$$
d\gamma=dn\,\delta(s)d^\times s\,dk,
$$ where $dn$ is a Haar measure on the unipotent group $N(\R)$,
$dk$ is a Haar measure on the compact group~$K$, $d^\times s$ is
given by
$$
d^\times s:=\prod_{i=1}^g\frac{ds_i}{s_i},
$$
and
\begin{equation}\label{eqhaarodd}
\source{squarefree-values-polynomial-discriminants}
\delta(s)=\prod_{k=1}^g s_k^{k^2-2kg};
\end{equation}
see \cite[(10.7)]{BG2}.

Since $s_i\geq c$ for every $i$, there exists a compact subset $N''$
of $N(\R)$ containing $(t)^{-1}N'\,(t)$ for all $t\in T'$. Since
$N''$, $K$, $\GG_1$ are compact and $R$ is bounded, the set $E=N''K\GG_1R$
is bounded. Then we have
\begin{equation}\label{eq:RE}
\source{squarefree-values-polynomial-discriminants}
N(\LL;M;X)\ll \int_{s_i\gg 1} \#\{B\in((s)\cdot XE)\cap\LL:|Q|(B)>M\}\delta(s)d^\times s.
\end{equation}




We denote the coordinates on $W$ by $b_{ij}$, for $1\leq i\leq j\leq
n$. These coordinates are eigenvectors for the action of $T$ on $W^*$,
the dual of $W$. Denote the $T$-weight of a coordinate $\alpha$ on
$W$, or more generally a product $\alpha$ of powers of such
coordinates, by $w(\alpha)$. An elementary computation shows that
\begin{equation}\label{wbij}
\source{squarefree-values-polynomial-discriminants}
w(b_{ij})=\left\{
\begin{array}{rcl}
t_i^{-1}t_j^{-1} &\mbox{ if }& i,j\leq g\\
t_i^{-1} &\mbox{ if }& i\leq g,\;j=g+1\\
t_i^{-1}t_{n-j+1} &\mbox{ if }& i\leq g,\; j>g+1\\
1 &\mbox{ if }& i=j=g+1\\
t_{n-j+1} &\mbox{ if }& i=g+1,\;j>g+1\\
t_{n-i+1}t_{n-j+1} &\mbox{ if }& i,j>g+1.
\end{array}
\right.
\end{equation}
Then the $(i,j)$-entry of any $B\in (s)\cdot XE$ is bounded by
$Xw(b_{ij})$, up to a multiplicative constant depending only on $\GG_1$
and $R$.

%For any field $k$ of characteristic not $2$, we say an element $B\in
%W(k)$ having nonzero discriminant is $k$-\emph{distinguished} if there
%exists a $g$-dimensional $k$-subspace isotropic with respect to the
%quadratic forms defined by the two symmetric matrices $A_0$ and
%$B$.
Let $W^\dist$ denote the subset of $W(\Q)$ consisting of
$\Q$-distinguished elements. Then $\LL$ is a subset of $W^\dist$. It
is shown in \cite[\S10.2]{BG2} that most of the lattice points in
$W^\dist$ lie inside the subspace $W_{00}$ consisting of symmetric
matrices $B$ whose $(i,j)$-entries are 0 whenever $i+j<n$. More
precisely, we have the following estimates from \cite[Propositions~10.5 and 10.7]{BG2}:

\begin{proposition}
\source{squarefree-values-polynomial-discriminants}
\notready\label{prop:bg2}
\source{squarefree-values-polynomial-discriminants}
We have
\begin{eqnarray}
% \nonumber % Remove numbering (before each equation)
  \int_{s_i\gg 1} \#\{B\in((s)\cdot XE)\cap({\textstyle\frac12} W(\Z)\setminus{\textstyle\frac12} W_{00}(\Z)):b_{11}=0\}\delta(s)d^\times s &=& O_\epsilon(X^{n(n+1)/2-1+\epsilon}) \\
  \label{eq:selbergpre}\int_{s_i\gg 1} \#\{B\in((s)\cdot XE)\cap\LL:b_{11}\neq 0\}\delta(s)d^\times s &=& o(X^{n(n+1)/2}).
\source{squarefree-values-polynomial-discriminants}
\end{eqnarray}
\end{proposition}

Therefore, to prove Theorem~\ref{thm:mainestimate}(b) when $n$ is odd, it remains to obtain a power-saving improvement of \eqref{eq:selbergpre} and estimate
\begin{equation}\label{eq:Qinv}
\source{squarefree-values-polynomial-discriminants}
\int_{s_i\gg 1} \#\{B\in ((s)\cdot XE)\cap \LL\cap
{\textstyle\frac12} W_{00}(\Z):|Q|(B)>M\}\delta(s)d^\times s.
\end{equation}


\subsection{Proof of Theorem~\ref{thm:mainestimate}(b)
  for odd $n$}\label{sgomodd}
\source{squarefree-values-polynomial-discriminants}

We begin with a power-saving improvement of \eqref{eq:selbergpre}.

\begin{proposition}
\source{squarefree-values-polynomial-discriminants}
\notready\label{propodd}
\source{squarefree-values-polynomial-discriminants}
  We have
  \begin{equation*}
\displaystyle\int_{s_i\gg 1} \#\{B\in ((s)\cdot XE)\cap W^\dist \cap {\textstyle\frac12} W(\Z): b_{11}\neq 0\}\delta(s)d^\times s =O_\epsilon(X^{n(n+1)/2-1/5+\epsilon}).
  \end{equation*}
\end{proposition}
\begin{proof}
We note that the $p$-adic density of elements in $W(\Z_p)$ that are
$\Q_p$-distinguished is bounded uniformly away from $1$. In fact, this
density is bounded above by $1 - \frac{g}{2g+1}+O(1/p)$ by 
\cite[\S10.7]{BG2}.  Then an application of the Selberg sieve exactly
as in \cite{ShTs} yields the result.
\end{proof}

We now estimate \eqref{eq:Qinv}. The $Q$-invariant can also be
given a weight by viewing the torus $T$ as sitting inside $G_0$ and
using \eqref{eq:weightG_0}. Namely,
\begin{equation}\label{eqQweight}
\source{squarefree-values-polynomial-discriminants}
w(Q)=\prod_{k=1}^gt_k^{-1}=\prod_{k=1}^gs_k^{-k}.
\end{equation}
Since the polynomial $Q$ is homogeneous of degree $g(g+1)/2$ in the
coefficients of $W_0$, we see that the $Q$-invariant of any $B\in
(s)\cdot XE$ is bounded by $X^{g(g+1)/2}w(Q)$, up to a multiplicative
constant depending only on $\GG_1$ and $R$.

Although we shall not use this fact directly, 
%it can be deduced, following the  
%proof of Proposition~\ref{prop:bg2} and using Propo\ref{propodd} that the
the points in $((s)\cdot XE)\cap \frac12W_{00}(\Z)$ with irreducible
invariant polynomial occur predominantly when the coordinates $s_i$
are so large that they force any (half-)integral point of $(s)\cdot
XE$ to lie inside $W_{00}$; this can be deduced by following the  
proof of Proposition~\ref{prop:bg2} and using Proposition~\ref{propodd}. 
On the other hand, since the weight of
the $Q$-invariant is a product of negative powers of $s_i$, the
$Q$-invariants of points in $((s)\cdot XE)\cap \frac12W_{00}(\Z)$
become large when the coordinates $s_i$ are small.  It is the tension
between these two requirements that underlies the proof of the following proposition, 
which gives the desired power-saving estimate for the
number of elements in $((s)\cdot XE) \cap \frac12W_{00}(\Z)$ having
large $|Q|$-invariant.

\begin{proposition}
\source{squarefree-values-polynomial-discriminants}
\notready\label{proplargeQbound}
\source{squarefree-values-polynomial-discriminants}
We have
\begin{equation*}
\int_{s_i\gg 1} \#\{B\in ((s)\cdot XE)\cap \LL\cap
{\textstyle\frac12} W_{00}(\Z):|Q|(B)>M\}\delta(s)d^\times s=O(\frac{1}{M}X^{n(n+1)/2}\log X).
\end{equation*}
\end{proposition}

\begin{proof}
First note that if an element in $\frac12 W_{00}(\Z)$ has
$(i,j)$-coordinate $0$ for some $i+j=n$, then the element has
discriminant $0$ and hence is not in $\LL$. Since the weight of $b_{i,n-i}$ is $s_i^{-1}$, to count points in $\LL$ it suffices to integrate only in the region where $s_i\ll X$ for all $i$.
Furthermore, the condition on the $|Q|$-invariant implies that
it suffices to integrate only in the region where $X^{g(g+1)/2}w(Q) \gg M$.



Let $S$ denote the set of coordinates of $W_{00}$, i.e.,
$S=\{b_{ij}:i+j\geq n\}$. For $(s)$ in the range $1\ll s_i\ll X$, we
have $Xw(\alpha)\gg 1$ for all $\alpha\in S$. Hence the number of lattice points in $(s)\cdot (XE)$ for $(s)$ in this range is $\ll \prod_{\alpha\in S}(Xw(\alpha))$. Therefore, 
\begin{eqnarray*}
&&\displaystyle\int_{s_i\gg 1} \#\{B\in((s)\cdot (XE))\cap \LL\cap {\textstyle\frac12} W_{00}(\Z):|Q(B)|>M\}\delta(s)d^\times s \\
&\ll&\displaystyle\int_{1\ll s_i\ll X,\,\,X^{g(g+1)/2}w(Q) \gg M} \prod_{\alpha\in S}\bigl(Xw(\alpha)\bigr)\delta(s)d^\times s \\
&\ll&\displaystyle\int_{1\ll s_i\ll X,\,\,X^{g(g+1)/2}w(Q) \gg M} X^{n(n+1)/2-g^2}\prod_{k=1}^gs_k^{2k-1}d^\times s \\
&\ll&\displaystyle\frac{1}{M}\int_{s_i=1}^X X^{n(n+1)/2-g^2+g(g+1)/2}w(Q)\prod_{k=1}^gs_k^{2k-1}d^\times s \\
&\ll&\displaystyle\frac{1}{M}\int_{s_i=1}^X X^{n(n+1)/2-g(g-1)/2}\prod_{k=1}^gs_k^{k-1}d^\times s \\
&\ll&\displaystyle\frac{1}{M}X^{n(n+1)/2}\log(X),
\end{eqnarray*}
where the second inequality follows from the definition
\eqref{eqhaarodd} of $\delta(s)$ and the computation \eqref{wbij} of
the weights of the coordinates $b_{ij}$, the third inequality follows
from the fact that $X^{g(g+1)/2}w(Q) \gg M$, the fourth inequality
follows from the computation of the weight of $Q$ in
\eqref{eqQweight}, and the $\log X$ factor comes from the integration
over $s_1$.
\end{proof}


The estimate in Theorem \ref{thm:mainestimate}(b) for odd $n$ now
follows from Theorem \ref{keymaporbit} and
Propositions~\ref{prop:bg2}, \ref{propodd} and~\ref{proplargeQbound},
in conjunction with the bound on the number of reducible polynomials
proved in Proposition~\ref{propredboundall}.
